% methods.tex
% CANDI: Confidence-Aware Neural Denoising Imputer for Epigenomic Data
\label{methods}
\subsection{Data}

We trained and evaluated CANDI using two datasets. We used data from the ENCODE Imputation Challenge (EIC) \cite{schreiber2023encode}, featuring 35 distinct assays measured across 50 biosamples, following the original EIC train-validation-test split; since CANDI uses read counts as input, we obtained the corresponding raw read files (BAM). To have a more extensive dataset, we systematically collected data for all biosamples in the ENCODE database that contained at least one experiment from the 35 assays of interest, yielding 3,064 biosamples that we merged using a cell type-based merging strategy into a final dataset of 361 merged cell types. For each biosample containing ChIP-seq experiments, we include the corresponding ChIP-seq control as an additional always-available input channel that is never masked, and we simulate varying sequencing depths by read downsampling (Downsampling Factor $\in \{1, 2, 4\}$). Full data collection and processing details are in Supplementary Methods.

\subsection{CANDI Architecture}

Let $A = 35$ be the number of assays and $L = 1{,}200$ the number of 25 bp positions in a $G = 30{,}000$ bp window. For each cell type, CANDI receives the epigenomic read-count matrix $\mathbf{M} \in \mathbb{R}^{N \times L \times (A+1)}$ (the $(A{+}1)$-th channel is the always-available ChIP-seq control), the one-hot DNA sequence $\mathbf{S} \in \{0,1\}^{N \times 4 \times G}$, and input and output experimental covariates $\mathbf{C}_{\text{in}}, \mathbf{C}_{\text{out}}$ (four per assay: $\log_2$ sequencing depth, run type, read length, and sequencing platform). Input counts are transformed using $\text{arcsinh}(x)$ before the encoder, which stabilizes variance across the dynamic range of read counts.

An encoder $\mathcal{E}$ maps $(\mathbf{M}, \mathbf{S}, \mathbf{C}_{\text{in}})$ to a latent representation $\mathbf{Z}$ through parallel convolution towers with per-layer MetadataCrossAttention that conditions features on covariates using FiLM (Feature-wise Linear Modulation) \cite{perez2018film}, a fusion layer, and $n_{\text{sab}} = 4$ transformer encoder layers ($n_{\text{head}} = 9$) with rotary positional embeddings (RoPE). For each assay $a$ and position $\ell$, three separate deconvolution decoders---each conditioned on the output covariates $\mathbf{C}_{\text{out}}$---predict the parameters of a negative binomial $(\hat{n}_{a,\ell}, \hat{p}_{a,\ell})$ over raw counts, a Gaussian $(\hat{\mu}_{a,\ell}, \hat{\sigma}^2_{a,\ell})$ over $\text{arcsinh}$-transformed signal $p$-values, and a peak probability $\text{peak}_{a,\ell} \in [0,1]$. The negative binomial captures the overdispersion of sequencing counts, and the predicted Gaussian variance captures position-varying (aleatoric) uncertainty (Supplementary Methods).

\subsection{Learning Objectives and Training Strategy}

Our approach leverages self-supervised learning, where the model learns to reconstruct deliberately corrupted versions of its own input, by combining three complementary strategies (Section \ref{results}): full assay masking (mimicking imputation), full loci masking (mimicking masked language modeling), and denoising via downsampling. We train CANDI by minimizing $\mathcal{L} = w_{\text{obs}}\, \mathcal{L}_{\text{obs}} + w_{\text{imp}}\, \mathcal{L}_{\text{imp}}$, where each of $\mathcal{L}_{\text{obs}}$ (unmasked positions) and $\mathcal{L}_{\text{imp}}$ (masked positions) is a weighted sum of the negative binomial (count), Gaussian (signal), and binary cross-entropy (peak) negative log-likelihoods, with default weights $w_{\text{count}} = 0.5$, $w_{\text{pval}} = 1.0$, $w_{\text{peak}} = 0.1$, $w_{\text{obs}} = 0.25$, and $w_{\text{imp}} = 1.0$. We trained CANDI on 5,000 selected, non-overlapping 30kb regions, excluding chromosome 21 and reserving it exclusively for testing (Supplementary Methods).

\subsection{Evaluation Metrics}

To evaluate imputation performance, we compute genome-wide Mean Squared Error (MSE), Pearson correlation, and Spearman correlation; all performance evaluations were made on chromosome 21. For peak prediction, we evaluate using precision, recall, and area under the ROC curve (AUROC) against ENCODE narrowPeak calls. To assess the reliability of the model's uncertainty estimates, we measure calibration---the fraction of observed values that fall within the model's predicted $C\%$ confidence interval---and the Concordance Index (C-index), which assesses whether the predicted distributions correctly rank-order the data while accounting for uncertainty (Supplementary Methods). As biological validation, we evaluate how well epigenomic features derived from the model's predictions---from observed, denoised, denoised+imputed, and latent ($\mathbf{Z}$) sources---predict gene expression levels (RNA-seq $\log(\text{TPM}+1)$) via nested cross-validation. We further assess whether CANDI improves the reproducibility of downstream chromatin state annotation: for each source we annotate the genome with an 18-state HMM (a soft-evidence HMM for CANDI's Gaussian predictions, so that predictive uncertainty propagates into the annotation) and quantify reproducibility with SAGAconf's robustness ratio~\cite{shahraki2024robust} (Supplementary Methods).
